\documentclass[11pt,letterpaper]{article}

\usepackage[spanish]{babel}
\usepackage[utf8]{inputenc}
\usepackage[T1]{fontenc}
\usepackage{geometry}
\geometry{margin=2.5cm}

\usepackage{setspace}
\onehalfspacing

\usepackage{hyperref}
\hypersetup{
    colorlinks=true,
    linkcolor=blue,
    urlcolor=blue
}

\usepackage{graphicx}
\usepackage{float}
\usepackage{booktabs}
\usepackage{array}
\usepackage{longtable}
\usepackage{amsmath, amssymb}
\usepackage{csquotes}
\usepackage{xcolor}
\usepackage{enumitem}
\usepackage{caption}
\usepackage{subcaption}


\usepackage{minted}

\setlength{\parskip}{6pt}
\setlength{\parindent}{0pt}


\newcommand{\universidad}{Universidad Católica de Temuco}
\newcommand{\facultad}{Facultad de Ingeniería}
\newcommand{\carrera}{Ingeniería Civil en Informática}
\newcommand{\asignatura}{Programación II}
\newcommand{\evaluacion}{Evaluación 3: Sistema de Gestión de Restaurante con ORM}
\newcommand{\equipo}{Gabriel Gutiérrez, Ailyn Melillán, Nicolás Cayumán}
\newcommand{\fechaentrega}{Noviembre 2025}



\title{\textbf{\evaluacion}}
\author{\equipo}
\date{\fechaentrega}

\begin{document}

\begin{titlepage}

   
    \begin{flushleft}
        \includegraphics[width=0.4\textwidth]{images/logouct.png}
    \end{flushleft}

    \vspace{-1cm} 

    \centering

    {\Large \universidad \par}
    \vspace{0.3cm}
    {\large \facultad \par}
    \vspace{0.3cm}
    {\large \carrera \par}

    \vspace{1.5cm}
    {\LARGE\bfseries \evaluacion \par}

    \vspace{1cm}
    {\large Asignatura: \asignatura \par}
    \vspace{0.5cm}
    {\large Docente: Guido Mellado \par}
    {\large Ayudantes: Fernando Valdés, Joaquín Cantero \par}

    \vfill
    {\large Integrantes:\par}
    {\large \equipo \par}

    \vfill
    {\large \fechaentrega \par}

\end{titlepage}

\tableofcontents
\newpage



\section{Introducción}

El presente documento expone el desarrollo del \textbf{Sistema de Gestión de Clientes, Pedidos, Menús e Ingredientes para un Restaurante}, correspondiente a la Evaluación 3.

Este sistema representa un avance significativo respecto a la Evaluación 2, agregando:

\begin{itemize}
    \item Persistencia profesional mediante \textbf{SQLAlchemy ORM}.
    \item Implementación completa de \textbf{CRUDs} modularizados.
    \item Manejo de \textbf{relaciones complejas} entre entidades (1:N y N:N).
    \item Interfaz gráfica profesional usando \textbf{CustomTkinter}.
    \item Programación funcional: \texttt{lambda}, \texttt{map}, \texttt{filter}, \texttt{reduce}.
    \item Generación de PDF: carta del menú y boleta del pedido.
    \item Gráficos estadísticos integrados con Matplotlib.
    \item Uso de patrones de diseño: \textbf{DTO} y \textbf{Facade}.
\end{itemize}

El objetivo del informe es describir detalladamente la arquitectura interna, el flujo del programa, las estructuras de datos, la interacción entre módulos y el cumplimiento de criterios de la pauta.


\section{Comparación entre Evaluación 2 y Evaluación 3}


\subsection{Mejoras más relevantes respecto a EV2}

La Evaluación 3 constituye una reingeniería total del proyecto. Las mejoras más importantes son:

\begin{itemize}
    \item \textbf{Migración completa a ORM}: desde clases en memoria (EV2) hacia un sistema persistente con SQLAlchemy (EV3).
    \item \textbf{Creación de CRUDs separados}: en EV2 toda la lógica estaba unida; en EV3 se modulariza según entidad.
    \item \textbf{Implementación de DTOs}: permiten desacoplar la interfaz gráfica del ORM.
    \item \textbf{Control real de stock}: se descuenta según recetas (N:N), no existía en EV2.
    \item \textbf{Interfaz moderna y multipestaña}: EV2 tenía una GUI mínima.
    \item \textbf{PDFs}: boleta y carta del restaurante.
    \item \textbf{Gráficos estadísticos}: inexistentes en EV2.
\end{itemize}

\section{Estructura General del Proyecto}

La estructura del proyecto es modular:

\begin{itemize}
    \item \textbf{app.py}: flujo principal del programa e interfaz gráfica.
    \item \textbf{database.py}: inicialización del ORM y la base de datos.
    \item \textbf{models.py}: modelos ORM (Cliente, Menú, Ingrediente, Pedido, Detalle, asociaciones).
    \item Carpeta \textbf{crud/}:
    \begin{itemize}
        \item \texttt{cliente\_crud.py}
        \item \texttt{crud\_pedido.py}
        \item \texttt{crud\_detallepedido.py}
        \item \texttt{ingrediente\_crud.py}
        \item \texttt{menu\_crud.py}
    \end{itemize}
    \item \textbf{BoletaFacade.py}: encapsula el proceso completo de generación de boletas.
    \item \textbf{ElementoMenu.py} y \textbf{Ingrediente.py}: DTOs.
    \item \textbf{graficos.py}: extracción de datos y construcción de gráficos.
\end{itemize}

Cada módulo cumple una única responsabilidad, siguiendo los principios SOLID.

\section{Modelo de Datos y ORM}

\subsection{Diagrama MER}

\begin{verbatim}
erDiagram
    CLIENTE ||--o{ PEDIDO : genera
    PEDIDO ||--o{ DETALLE_PEDIDO : contiene
    MENU ||--o{ DETALLE_PEDIDO : vendido
    MENU ||--o{ MENU_INGREDIENTE : receta
    INGREDIENTE ||--o{ MENU_INGREDIENTE : requerido
\end{verbatim}

\subsection{Definición de modelos ORM}

Ejemplo de un modelo utilizado:

\begin{minted}[fontsize=\small]{python}
class MenuIngrediente(Base):
    __tablename__ = "menu_ingredientes"

    id = Column(Integer, primary_key=True)
    menu_id = Column(Integer, ForeignKey("menus.id"))
    ingrediente_id = Column(Integer, ForeignKey("ingredientes.id"))
    cantidad_requerida = Column(Float)

    menu = relationship("Menu", back_populates="ingredientes_asociados")
    ingrediente = relationship("Ingrediente")
\end{minted}

Este diseño permite controlar el stock de cada menú de forma precisa.

\section{CRUDs (Create, Read, Update, Delete)}

\subsection{Ejemplo: CRUD de Menús}

\begin{minted}[fontsize=\small]{python}
def create_menu(self, db, nombre, precio, icono):
    nuevo = Menu(nombre=nombre, precio=precio, icono=icono)
    db.add(nuevo)
    db.commit()
    db.refresh(nuevo)
    return nuevo
\end{minted}

\textbf{Aspectos importantes:}

\begin{itemize}
    \item \texttt{db.add() + db.commit()} permite persistir cambios.
    \item \texttt{db.refresh()} obtiene valores actualizados del ORM.
    \item Todas las operaciones CRUD están aisladas en un módulo específico.
\end{itemize}

% =============================================================
\section{Interfaz Gráfica (CustomTkinter)}
% =============================================================

El archivo \texttt{app.py} organiza todas las vistas en formato de pestañas:

\begin{minted}[fontsize=\small]{python}
self.tabview.add("Pedido")
self.tabview.add("Ingredientes")
self.tabview.add("Menús")
self.tabview.add("Clientes")
self.tabview.add("Gráficos")
self.tabview.add("Boleta")
\end{minted}

Cada pestaña interactúa con los CRUDs y DTOs.

% =============================================================
\section{Programación Funcional}
% =============================================================

\subsection{Uso de \texttt{map}}

\begin{minted}[fontsize=\small]{python}
nombres_clientes = list(map(
    lambda c: f"{c.id} - {c.nombre}",
    clientes
))
\end{minted}

\subsection{Uso de \texttt{filter}}

\begin{minted}[fontsize=\small]{python}
menus_disponibles = list(filter(
    lambda m: self.menu_crud.esta_disponible(self.db_session, m),
    menus_db
))
\end{minted}

\subsection{Uso de \texttt{reduce}}

\begin{minted}[fontsize=\small]{python}
total = reduce(
    lambda subtotal, item: subtotal + item.precio * item.cantidad,
    self.pedido.menus, 0.0
)
\end{minted}

\subsection{Uso de \texttt{lambda} en eventos}

\begin{minted}[fontsize=\small]{python}
tarjeta.bind("<Enter>", lambda e: tarjeta.configure(border_color="red"))
\end{minted}

% =============================================================
\section{Generación de Boleta (Patrón Facade)}
% =============================================================

\subsection{Arquitectura del Facade}

\begin{minted}[fontsize=\small]{python}
class BoletaFacade:
    def procesar_pedido(self, pedido_dto):
        pedido_orm = self.pedido_crud.crear_pedido(self.db, pedido_dto)
        self.detalle_crud.crear_detalles(self.db, pedido_orm, pedido_dto)
        self.actualizar_stock(pedido_orm)
        return self.generar_pdf(pedido_orm)
\end{minted}

\textbf{Responsabilidades encapsuladas:}

\begin{itemize}
    \item Registrar el pedido.
    \item Registrar detalles.
    \item Descontar stock.
    \item Generar boleta en PDF.
\end{itemize}

% =============================================================
\section{Flujo Completo del Sistema}
% =============================================================

\textbf{1. Inicialización}

\begin{minted}[fontsize=\small]{python}
create_db_and_tables()
self.db = SessionLocal()
\end{minted}

\textbf{2. Construcción de pestañas y DTOs}

\begin{minted}[fontsize=\small]{python}
self.pedido = Pedido()
\end{minted}

\textbf{3. Agregar menú al carrito}

\begin{minted}[fontsize=\small]{python}
self.pedido.menus.append(menu_dto)
\end{minted}

\textbf{4. Validación de stock}

\begin{minted}[fontsize=\small]{python}
if not self.menu_crud.esta_disponible(self.db, menu_orm):
    return
\end{minted}

\textbf{5. Cálculo del total}

\begin{minted}[fontsize=\small]{python}
total = reduce(lambda t, m: t + m.precio*m.cantidad, self.pedido.menus, 0)
\end{minted}

\textbf{6. Persistencia del pedido vía Facade}

\begin{minted}[fontsize=\small]{python}
ruta_pdf = facade.procesar_pedido(self.pedido)
\end{minted}

\textbf{7. Visualización de la boleta en PDF}

% =============================================================
\section{Conclusiones}
% =============================================================

El sistema final cumple todos los requisitos establecidos por la pauta de la Evaluación 3, integrando:

\begin{itemize}
    \item Interfaz gráfica robusta.
    \item ORM profesional con SQLAlchemy.
    \item CRUDs completos y modularizados.
    \item Programación funcional aplicada.
    \item Flujo completo del proceso de compra.
    \item Generación de PDFs.
    \item Estadísticas gráficas.
\end{itemize}

El avance respecto a la Evaluación 2 es evidente tanto en arquitectura como en funcionalidad y calidad del código.

\end{document}







